% ============================================================================
\clearpage
\section*{v4.9.0-r55 Source-Only: A Readout per Particle Kind from the Subgroups of \texorpdfstring{\(A_4\)}{A4}}
\addcontentsline{toc}{section}{v4.9.0-r55 Source-Only: A Readout per Particle Kind from the Subgroups of A4}
\label{sec:v490-r55-residual-readouts}
% ============================================================================

\paragraph{Purpose.}
Mixing measures the difference between the bases in which two particle kinds read their masses.  The author asked what follows if each particle
kind needs its own readout.  This layer records the pre-registered test H39 (note \texttt{llm/Vorab\_Notiz\_H39\_Restsymmetrie\_Lesarten.md}, commit
\texttt{5c916b8}, fixed before computing), with readouts taken from the subgroups of the common symmetry \(A_4\) of the Dixon store and the triality
form (v4.9.0-r33).  Earlier layers are not rewritten.

\noindent\textbf{Source-only status.}
This v4.9.0-r55 layer is source-only.  It does not build a PDF, does not assign a new DOI,
and does not authorize a public release.  The v4.9.0 Version DOI
\href{https://doi.org/10.5281/zenodo.23534954}{10.5281/zenodo.23534954} and the Concept DOI
\href{https://doi.org/10.5281/zenodo.19210728}{10.5281/zenodo.19210728} are unchanged by it.

\subsection*{v4.9.0-r55 rule}
\label{subsec:v490-r55-rule}

A readout of a particle kind is the subgroup \(G\subset A_4\) under which its mass coupling is invariant; it fixes its eigenbasis wholly or in part.
The octonion sectors do not distinguish \(u\) from \(d\) or \(\nu\) from \(e\) (v4.9.0-r49).  What the structure does distinguish is v4.9.0-r44:
Dirac couplings (\(u\), \(d\), \(e\)) are circulant, hence invariant under \(G_D=\langle\tau\rangle\cong\mathbb Z_3\); neutrinos admit a Majorana term and read
with another subgroup \(G_M\), which the structure does not fix.  All subgroups except \(\langle\tau\rangle\) and \(A_4\) were computed.  The note
corrects an announcement in the discussion: tribimaximal mixing needs \(S_4\), not \(A_4\).

\subsection*{v4.9.0-r55 results}
\label{subsec:v490-r55-results}

\paragraph{Theoremlet v4.9.0-r55.1.}
\begin{enumerate}[leftmargin=2em]
\item[(a)] \(A_4\) has subgroups of orders \(1,2,3,4,12\) with multiplicities \(1,3,4,1,1\).  With \(G_M\) the Klein group all \(|U_{ij}|^2=\tfrac13\); with
\(G_M\) one of the three other \(\mathbb Z_3\) the basis is fixed; neither reaches the NuFIT \(3\sigma\) ranges.
\item[(b)] With \(G_M\cong\mathbb Z_2\) one column is fixed and one angle is free.  The \(3\sigma\) ranges are reached exactly when the fixed column is \(\nu_2\)
(trimaximal mixing TM2): then \(\sin^2\theta_{12}=1/(3\cos^2\theta_{13})\in[0.3402,0.3414]\), while \(\sin^2\theta_{23}\) ranges over \([0.399,0.601]\).
\item[(c)] With the same readout for \(u\) and \(d\) the CKM matrix is the identity at leading order.
\end{enumerate}

\paragraph{Proof sketch.}
Subgroups by closure, joint eigenspaces of each subgroup, a \(U(2)\) rotation in a two-dimensional eigenspace, all column assignments (script
\texttt{llm/k95\_test\_h39.py}).  The split of (b) by assignment is a check after the evaluation, appended after the first run
(\texttt{k95\_out\_erstlauf.txt}), whose criterion values agree.  \(\square\)

\paragraph{Verdict and reading.}
H39a, b, c are confirmed.  Over all assignments the remaining quantities spread wider than the pre-registered widths, so by the wording this is a
fit with one angle.
\begin{enumerate}[leftmargin=2em]
\item The readout rule works qualitatively: Dirac sectors with \(\mathbb Z_3\) and neutrinos with \(\mathbb Z_2\) from the same \(A_4\) give no quark mixing
at leading order and large lepton mixing.  This is the difference between CKM and PMNS that the tests of v4.9.0-r50 to r54 did not produce.
\item Quantitatively it is the known TM2 pattern of the \(A_4\) literature; the structure supplies the group and the rule, not a new pattern.  It fixes
\(\sin^2\theta_{12}\approx0.341\), inside the NuFIT \(3\sigma\) range.  A comparison with newer data, not part of the note: JUNO measures
\(\sin^2\theta_{12}=0.3092\pm0.0087\), \(3.6\sigma\) below.  \(\theta_{23}\) stays free.
\item Corrections, for instance from the mask of v4.9.0-r47, could lower \(\theta_{12}\) and produce the small CKM angles; that would be a test of its own.
\end{enumerate}

\subsection*{v4.9.0-r55 status in the three frames}
\label{subsec:v490-r55-frames}

\begin{center}
\small
\begin{tabular}{@{}>{\raggedright\arraybackslash}p{0.30\textwidth}>{\raggedright\arraybackslash}p{0.62\textwidth}@{}}
\toprule
result & M / T / P\\
\midrule
r55.1(a)--(c) & statements on the slot symmetry \(A_4\), which acts on the slots and commutes with the frame maps on the \(\mathbb C\) factor: the same in M and T; P not evaluated\\
\bottomrule
\end{tabular}
\end{center}

\subsection*{v4.9.0-r55 register entries}
\label{subsec:v490-r55-register}

\begin{center}
\small
\begin{tabular}{@{}>{\raggedright\arraybackslash}p{0.34\textwidth}>{\raggedright\arraybackslash}p{0.58\textwidth}@{}}
\toprule
statement & status after v4.9.0-r55\\
\midrule
readout per particle kind (Dirac \(\mathbb Z_3\), Majorana \(\mathbb Z_2\)) & qualitatively right ordering of CKM and PMNS; quantitatively TM2\\
\(\sin^2\theta_{12}=1/(3\cos^2\theta_{13})\approx0.341\) & sharp; inside NuFIT \(3\sigma\), \(3.6\sigma\) above JUNO 2025\\
\(\theta_{23}\), CKM angles & not fixed at leading order; corrections open\\
\bottomrule
\end{tabular}
\end{center}

\subsection*{v4.9.0-r55 ledger and guard update}
\label{subsec:v490-r55-ledger-guard}

\[
R_{\rm rr}^{v490\text{-}r55}=(1,0,0)
\]
for (readout rule with qualitatively right ordering; quantitative mixing not explained; corrections open).

The layer spends no guard:
\[
R_{\rm glob}^{(v490\text{-}r55,+)}=(0,1,1,1,1),
\qquad
\Gamma_{\rm guard}^{v490\text{-}r55}=(1,1,1,1,1),
\qquad
\beta_4=1/256\ \text{unspent}.
\]
It does not imply exact finite \(C5\) zero closure, exact all-readout terminality,
\(C6/T2\), a global photon-clock/RH theorem, or global all-readout closure.
